\begin{exsol}{4}{Cabral 3.9 Car parts}{62}{Foundations}{A2}
\asked Of two parts of a car's taillight, the plastic exterior cover and the light bulb, say
which a carmaker is more likely to make in-house, and why.

\method The make-or-buy trade-off of \S3.4 (pp.~54--56): vertical boundaries balance
\emph{investment incentives} (specific assets and the hold-up problem) against
\emph{performance incentives} (an outside supplier who can be replaced works harder), with
intermediate forms such as long-term supplier relations (p.~55). The mould for a car part is
Cabral's own example of a specific asset (\S13.1, p.~338). Course cases: GM's integration and
Toyota's outside suppliers (\MR{} Ch.~1, pp.~4--6).

\solution
\subparagraph{Compare the two parts on the dimensions that decide make or buy.}
\begin{center}\small
\begin{tabularx}{\linewidth}{@{}>{\raggedright\arraybackslash}p{3.3cm}XX@{}}
\toprule
 & Exterior cover & Light bulb \\
\midrule
Asset specificity & high: shape, colour and fit are unique to one model; the mould is worthless
for anything else & low: a standard part sold to many carmakers and in the aftermarket \\
Number of suppliers ex post & one (whoever holds the mould) & many \\
Scale economies & limited to one model's volume, which the carmaker can reach itself & large:
a specialist pools demand from the whole industry \\
Design coordination & tight: must match the body, the lamp housing and the styling & minimal:
fits a standard socket \\
Supplier know-how & moulding plastics, close to what a carmaker already does & specialized
lighting technology \\
\bottomrule
\end{tabularx}
\end{center}

\subparagraph{The bulb is bought.} It is a standard part: no investment is specific to one
buyer, so no hold-up arises; competing suppliers keep prices near cost and the threat of
switching keeps performance high (p.~55); and a specialist that serves the whole industry
reaches scale economies that no single carmaker could. Making it in-house would give up the
market's performance incentive and its scale for no gain.

\subparagraph{The cover is the candidate for making in-house.} Producing it requires a
model-specific mould. Once an outside supplier has sunk that investment, the carmaker can
squeeze the price (the supplier has no other customer for the mould); once the carmaker has
committed to the supplier's mould in the middle of a model run, the supplier can raise the price
(no one else can supply the part at short notice). Either side can hold up the other, which is
Cabral's coal-mine and power-plant logic (p.~55) and the mould example of \S13.1 (p.~338).
Ownership removes the bargaining between the two sides. Close design coordination with the body
and the small extra know-how needed point in the same direction.

\subparagraph{Qualification: intermediate forms.} ``More likely in-house'' is a comparative
statement. The hold-up problem can also be handled without full ownership: a long-term
relation with a few trusted suppliers who are willing to invest in buyer-specific assets (the
Japanese supplier system, p.~55), or an arrangement where the carmaker owns the specific mould
while the supplier runs it, which takes the specific asset off the supplier's books and lets the
carmaker move the mould to another supplier if the price is raised. This is the
GM-versus-Toyota contrast of Aula~1: GM's volumes let it run efficient parts plants and it
integrated (\MR{} p.~5), while Toyota, too small to reach efficient scale in parts, used few
outside suppliers in long-term relations for complex components (pp.~4--6).

\interp The make-or-buy answer depends on the transaction, not on the industry: the same
carmaker should buy the commodity and make, or govern closely, the specific part. Specificity
creates hold-up, which favours integration; standardization makes market discipline available,
which favours buying. Item~30 (Cabral 13.10) puts numbers on the hold-up side.

\verify Limit cases. If the cover were a standard shape used across many models and makers, its
mould would stop being specific and the cover would move to the ``buy'' column. If a bulb were a
custom design for one model (a model-specific LED unit, say), it would move towards the ``make or
govern closely'' column. The prediction follows the specificity of the asset, not the physical
part.

\answer{The plastic exterior cover is more likely made in-house: its model-specific mould is a
specific asset, so outsourcing exposes both sides to hold-up, and design coordination with the
body is tight. The bulb is a standard part with many suppliers and large scale economies, so it
is bought and market competition disciplines the supplier (\S3.4, p.~55). Long-term supplier
relations (Toyota) are the intermediate solution.}
\end{exsol}
